% !TITLE 黄海舟中日人索句并见日俄战争地图
% !AUTHOR 秋瑾
% !DYNASTY 近代
% !GENRE 七言律诗
% !ORDER 3

\begin{poem}
万里乘风\textsuperscript{1}去复来，只身东海挟春雷\textsuperscript{2}。
忍看图画移颜色\textsuperscript{3}，肯使江山付劫灰\textsuperscript{4}。
浊酒难销忧国泪，救时应仗出群才\textsuperscript{5}。
拼将十万头颅血，须把乾坤\textsuperscript{6}力挽回。
\end{poem}

\begin{annotations}
\notetext{1}{乘风：驾着风，比喻志向远大。语出“乘风破浪”。秋瑾曾两次东渡日本。}
\notetext{2}{挟春雷：带着春天的雷声。比喻带着革命的决心与力量。}
\notetext{3}{图画移颜色：地图改变了颜色。指日俄战争中俄国将在中国东北的权益转让给日本——中国的领土在地图上被换了颜色，却不能自己做主。}
\notetext{4}{劫灰：劫火之后留下的灰烬。劫，佛教中指毁灭一切的灾难。怎能让祖国山河化为劫后的灰烬？}
\notetext{5}{救时应仗出群才：拯救时局应当依靠那些出类拔萃的人才。仗，依靠；出群才，超群的人才。}
\notetext{6}{乾坤：天地，代指国家、天下大局。}
\end{annotations}
